\setcounter{ex}{0}

% --- CÂU 1 ---
\begin{ex}
	Cho $\tan\alpha = -\dfrac{3}{4}$ với $\dfrac{\pi}{2} < \alpha < \pi$. Giá trị của $\sin\alpha$ bằng
	\choice
	{$\dfrac{4}{5}$}
	{$-\dfrac{4}{5}$}
	{$-\dfrac{3}{5}$}
	{\True $\dfrac{3}{5}$}
	\loigiai{
		Vì $\dfrac{\pi}{2} < \alpha < \pi$ nên $\sin\alpha > 0$ và $\cos\alpha < 0$. Ta có:
		$$\dfrac{1}{\cos^2\alpha} = 1 + \tan^2\alpha = 1 + \dfrac{9}{16} = \dfrac{25}{16} \implies \cos\alpha = -\dfrac{4}{5}.$$
		Khi đó $\sin\alpha = \tan\alpha \cdot \cos\alpha = \left(-\dfrac{3}{4}\right) \cdot \left(-\dfrac{4}{5}\right) = \dfrac{3}{5}$.
	}
\end{ex}

% --- CÂU 2 ---
\begin{ex}
	Cho $\sin x + \cos x = \dfrac{1}{3}$. Giá trị của $\sin 2x$ bằng
	\choice
	{\True $-\dfrac{8}{9}$}
	{$\dfrac{8}{9}$}
	{$\dfrac{4}{9}$}
	{$-\dfrac{4}{9}$}
	\loigiai{
		Bình phương hai vế của đẳng thức $\sin x + \cos x = \dfrac{1}{3}$, ta được:
		$$(\sin x + \cos x)^2 = \dfrac{1}{9} \iff \sin^2 x + \cos^2 x + 2\sin x\cos x = \dfrac{1}{9} \iff 1 + \sin 2x = \dfrac{1}{9} \iff \sin 2x = -\dfrac{8}{9}.$$
	}
\end{ex}

% --- CÂU 3 ---
\begin{ex}
	Tập xác định của hàm số $y = \dfrac{\tan x}{\cos 2x}$ là
	\choice
	{$\mathbb{R} \setminus \left\{\dfrac{\pi}{4} + \dfrac{k\pi}{2},\ k \in \mathbb{Z}\right\}$}
	{$\mathbb{R} \setminus \left\{\dfrac{k\pi}{2},\ k \in \mathbb{Z}\right\}$}
	{\True $\mathbb{R} \setminus \left\{\dfrac{\pi}{2} + k\pi;\ \dfrac{\pi}{4} + \dfrac{k\pi}{2},\ k \in \mathbb{Z}\right\}$}
	{$\mathbb{R} \setminus \left\{\dfrac{\pi}{2} + k\pi,\ k \in \mathbb{Z}\right\}$}
	\loigiai{
		Hàm số xác định khi và chỉ khi:
		$$\heva{\cos x &\ne 0 \\ \cos 2x &\ne 0} \iff \heva{x &\ne \dfrac{\pi}{2} + k\pi \\ 2x &\ne \dfrac{\pi}{2} + k\pi} \iff \heva{x &\ne \dfrac{\pi}{2} + k\pi \\ x &\ne \dfrac{\pi}{4} + \dfrac{k\pi}{2}}\ (k \in \mathbb{Z}).$$
		Vậy tập xác định của hàm số là $\mathscr{D} = \mathbb{R} \setminus \left\{\dfrac{\pi}{2} + k\pi;\ \dfrac{\pi}{4} + \dfrac{k\pi}{2},\ k \in \mathbb{Z}\right\}$.
	}
\end{ex}

% --- CÂU 4 ---
\begin{ex}
	Hàm số nào dưới đây là hàm số lẻ?
	\choice
	{$y = |\sin x|$}
	{\True $y = \sin x\cos x$}
	{$y = x^2 + \cos x$}
	{$y = \cos x + 1$}
	\loigiai{
		Xét hàm số $f(x) = \sin x\cos x$ có tập xác định $\mathscr{D} = \mathbb{R}$ là tập đối xứng.\\
		Với mọi $x \in \mathbb{R}$, ta có $f(-x) = \sin(-x)\cos(-x) = -\sin x\cos x = -f(x)$ nên $y = \sin x\cos x$ là hàm số lẻ.
	}
\end{ex}

% --- CÂU 5 ---
\begin{ex}
	Nghiệm của phương trình $\cos\left(x - \dfrac{\pi}{4}\right) = -\dfrac{\sqrt{2}}{2}$ là
	\choice
	{$x = \pi + k\pi\ (k \in \mathbb{Z})$}
	{$x = \dfrac{\pi}{4} \pm \dfrac{\pi}{4} + k2\pi\ (k \in \mathbb{Z})$}
	{$x = \pm\dfrac{3\pi}{4} + k2\pi\ (k \in \mathbb{Z})$}
	{\True $x = \pi + k2\pi$ hoặc $x = -\dfrac{\pi}{2} + k2\pi\ (k \in \mathbb{Z})$}
	\loigiai{
		Ta có:
		$$\cos\left(x - \dfrac{\pi}{4}\right) = -\dfrac{\sqrt{2}}{2} \iff \cos\left(x - \dfrac{\pi}{4}\right) = \cos\dfrac{3\pi}{4} \iff \hoac{x - \dfrac{\pi}{4} &= \dfrac{3\pi}{4} + k2\pi \\ x - \dfrac{\pi}{4} &= -\dfrac{3\pi}{4} + k2\pi} \iff \hoac{x &= \pi + k2\pi \\ x &= -\dfrac{\pi}{2} + k2\pi}\ (k \in \mathbb{Z}).$$
	}
\end{ex}

% --- CÂU 6 ---
\begin{ex}
	Số nghiệm của phương trình $\tan x = 1$ trên đoạn $[-\pi; 2\pi]$ là
	\choice
	{$2$}
	{$6$}
	{$4$}
	{\True $3$}
	\loigiai{
		Ta có $\tan x = 1 \iff x = \dfrac{\pi}{4} + k\pi\ (k \in \mathbb{Z})$.\\
		Vì $x \in [-\pi; 2\pi]$ nên $-\pi \le \dfrac{\pi}{4} + k\pi \le 2\pi \iff -\dfrac{5}{4} \le k \le \dfrac{7}{4}$.\\
		Do $k \in \mathbb{Z}$ nên $k \in \{-1; 0; 1\}$. Vậy phương trình có $3$ nghiệm trên đoạn $[-\pi; 2\pi]$.
	}
\end{ex}

% --- CÂU 7 ---
\begin{ex}
	Giá trị nhỏ nhất của hàm số $y = 3 - 2\sin^2 x$ bằng
	\choice
	{$3$}
	{$5$}
	{$-1$}
	{\True $1$}
	\loigiai{
		Vì $0 \le \sin^2 x \le 1, \forall x \in \mathbb{R}$ nên $1 \le 3 - 2\sin^2 x \le 3$.\\
		Vậy giá trị nhỏ nhất của hàm số bằng $1$ (khi $\sin^2 x = 1$).
	}
\end{ex}

% --- CÂU 8 ---
\begin{ex}
	Tìm $x$ để ba số $x - 1,\ x + 3,\ 2x + 3$ theo thứ tự lập thành một cấp số cộng.
	\choice
	{$x = 6$}
	{$x = 5$}
	{$x = 3$}
	{\True $x = 4$}
	\loigiai{
		Ba số $x - 1,\ x + 3,\ 2x + 3$ theo thứ tự lập thành một cấp số cộng khi và chỉ khi:
		$$(x - 1) + (2x + 3) = 2(x + 3) \iff 3x + 2 = 2x + 6 \iff x = 4.$$
	}
\end{ex}

% --- CÂU 9 ---
\begin{ex}
	Cho cấp số cộng $(u_n)$ có $u_1 = -4$ và công sai $d = 3$. Tổng $20$ số hạng đầu bằng
	\choice
	{\True $490$}
	{$500$}
	{$480$}
	{$470$}
	\loigiai{
		Tổng $20$ số hạng đầu của cấp số cộng là:
		$$S_{20} = 20u_1 + \dfrac{20 \cdot 19}{2}d = 20 \cdot (-4) + 190 \cdot 3 = -80 + 570 = 490.$$
	}
\end{ex}

% --- CÂU 10 ---
\begin{ex}
	Cho hình chóp $S.ABCD$ có đáy $ABCD$ là hình bình hành tâm $O$, $M, N$ lần lượt là trung điểm của $SC, SD$ (hình vẽ dùng chung cho câu 10, 11, 12). Giao tuyến của hai mặt phẳng $(SAC)$ và $(SBD)$ là
	\begin{center}
		\begin{tikzpicture}[scale=0.85, font=\footnotesize, line join=round, line cap=round]
			\coordinate (A) at (0,0);
			\coordinate (B) at (3.2,0);
			\coordinate (D) at (1.3,1.1);
			\coordinate (C) at (4.5,1.1);
			\coordinate (O) at ($(A)!0.5!(C)$);
			\coordinate (S) at (2.2,3.4);
			\coordinate (M) at ($(S)!0.5!(C)$);
			\coordinate (N) at ($(S)!0.5!(D)$);
			\coordinate (I) at (intersection of A--M and S--O);
			\draw[thick] (S)--(A) (S)--(B) (S)--(C) (A)--(B)--(C);
			\draw[dashed] (S)--(D) (A)--(D)--(C) (A)--(C) (B)--(D) (S)--(O) (A)--(M) (M)--(N);
			\foreach \p/\pos in {S/above, A/below left, B/below, C/right, D/left, O/below, M/above right, N/above left, I/left}
				\node[\pos] at (\p) {$\p$};
			\foreach \p in {S,A,B,C,D,O,M,N,I}
				\fill (\p) circle (1.2pt);
		\end{tikzpicture}
	\end{center}
	\choice
	{$AC$}
	{$SA$}
	{\True $SO$}
	{$BD$}
	\loigiai{
		Ta có $S \in (SAC) \cap (SBD)$ và $O = AC \cap BD \in (SAC) \cap (SBD)$ nên $(SAC) \cap (SBD) = SO$.
	}
\end{ex}

% --- CÂU 11 ---
\begin{ex}
	Giao điểm của đường thẳng $AM$ và mặt phẳng $(SBD)$ là giao điểm của $AM$ với
	\choice
	{$SB$}
	{$BD$}
	{$SD$}
	{\True $SO$}
	\loigiai{
		Vì $AM \subset (SAC)$ và $(SAC) \cap (SBD) = SO$ nên giao điểm của $AM$ và $(SBD)$ là giao điểm của $AM$ với $SO$.
	}
\end{ex}

% --- CÂU 12 ---
\begin{ex}
	Khẳng định nào sau đây đúng?
	\choice
	{$MN \parallel (SAD)$}
	{$MN \parallel (SBC)$}
	{$MN \parallel (SBD)$}
	{\True $MN \parallel (ABCD)$}
	\loigiai{
		Vì $M, N$ lần lượt là trung điểm của $SC, SD$ nên $MN$ là đường trung bình của tam giác $SCD \implies MN \parallel CD$.\\
		Mà $CD \subset (ABCD)$ và $MN \not\subset (ABCD)$ nên $MN \parallel (ABCD)$.
	}
\end{ex}
